\documentclass[12pt]{article}
\usepackage{notes-common}
\geometry{margin=0.35in}
\AtBeginDocument{\fontsize{14}{17}\selectfont}

% Complete / filled-in reference version of electric-pe-notes.tex (the
% skeleton students fill in by hand). Same structure and wording, but
% every blank carries the answer, and the student-drawn prompts (ball
% above ground + hole, charge in a field) get an actual figure.
%
% Restructured 2026-09-24 per Seth: gravity (height, then a hole) first,
% then a uniform E field for BOTH + and - charges, then the full
% (non-uniform) Newtonian form of gravitational PE with its own graph
% ("higher is more PE" even though U<0 everywhere), then electric PE of
% point charges (with a sketch of U vs. r right after U(q_1,q_2), for
% both opposite and like charges), and finally the dipole.
\newcommand{\ANS}[1]{\textcolor{accent}{#1}}

\begin{document}

\noteshead{Electric Potential Energy \textcolor{accent}{(completed)}}{}

\goal{
\begin{itemize}[itemsep=5pt,parsep=2pt,topsep=2pt]
\item Build the idea of potential energy with gravity first: height above
the ground, and what happens if you keep going \emph{below} it.
\item See the same idea for a charge pushed against a uniform field ---
for both a $+$ and a $-$ charge.
\item Write the \emph{full} (non-uniform) form of gravitational PE, and see
that ``higher'' always means ``more PE,'' even though the value is negative.
\item Use that same point-source structure to build the electric PE of two
point charges, then a whole system of them.
\item Combine two point charges' PE to get the potential energy of a
dipole in a field.
\end{itemize}}

%============================================================
\nsec{Gravitational PE: height, then a hole}

A ball sits a height $h$ above the ground. Ground level is our reference:
$U=0$ there.

\begin{center}
\begin{tikzpicture}[x=1cm,y=1cm]
  \draw[brown!60!black, line width=1.2pt] (-3,0) -- (-0.9,0);
  \draw[brown!60!black, line width=1.2pt] (0.9,0) -- (3,0);
  \foreach \x in {-2.7,-2.1,-1.5}{\draw[brown!60!black] (\x,0) -- (\x-0.2,-0.25);}
  \foreach \x in {1.4,2.0,2.6}{\draw[brown!60!black] (\x,0) -- (\x-0.2,-0.25);}
  \draw[brown!60!black, line width=1.2pt] (-0.9,0) -- (-0.9,-2) -- (0.9,-2) -- (0.9,0);
  \filldraw[gray!30, opacity=0.6] (0,1.7) circle (0.35);
  \draw[gray!60!black, opacity=0.6] (0,1.7) circle (0.35);
  \node[scale=0.75, accent] at (1.05,1.85) {start};
  \filldraw[gray!30] (0,-1.65) circle (0.35);
  \draw[gray!60!black] (0,-1.65) circle (0.35);
  \node[scale=0.75, accent] at (0.95,-1.65) {end};
  \draw[accent, dashed, ->] (-2.3,0) -- (-2.3,1.7);
  \node[accent, scale=0.8] at (-2.65,0.85) {$h$};
  \draw[accent, dashed, ->] (2.3,0) -- (2.3,-2);
  \node[accent, scale=0.8] at (2.65,-1) {$d$};
  \node[scale=0.75, blankink] at (-1.9,-1.15) {ground};
  \node[scale=0.75] at (0,0.28) {$h=0$};
\end{tikzpicture}
\end{center}

\[
U(h) = \ANS{mgh}
\]

Now let the ball fall past ground level, into a hole of depth $d$.

\[
U(\text{bottom}) = \ANS{-mgd}
\]

{\small\color{accent} Ground is still where $h=0$. The bottom of the hole is
\emph{below} that reference, at height $-d$, so $U=mg(-d)=-mgd$ --- our
first negative PE.}

\[
\Delta U = U(\text{bottom}) - U(h) = \ANS{-mgd - mgh \;=\; -mg(h+d)}
\]

\checkpoint{The ball's PE at the bottom of the hole came out negative. Is
that a problem?}

{\small\color{accent} No --- ``zero'' was just our choice of ground level.
PE can absolutely be negative below wherever you chose to measure from. Keep
this in mind: we'll see the exact same thing with opposite charges.}

%============================================================
\nsec{The same idea, for a charge in a uniform field}

Replace gravity with a uniform electric field $\vec E$ pointing
\emph{down} (same visual sense as gravity), and try it with both signs of
charge.

\begin{center}
\begin{tikzpicture}[x=1cm,y=1cm]
  % legend, top right, clear of everything else
  \draw[cyan!65!black, ->, line width=0.8pt] (7.6,1.9) -- (7.6,1.0);
  \node[cyan!65!black, scale=0.8] at (8.05,1.45) {$\vec E$};
  % --- left panel: +q ---
  \draw[cyan!65!black, ->, line width=0.8pt] (-6.6,1.9) -- (-6.6,-1.9);
  \draw[cyan!65!black, ->, line width=0.8pt] (-2.6,1.9) -- (-2.6,-1.9);
  \filldraw[red!70!black] (-4.6,1.15) circle (0.1);
  \node[red!70!black, scale=0.85] at (-4.6,1.55) {$+q$};
  \draw[accent, ->, line width=0.9pt] (-4.95,0.95) -- (-4.95,0.05);
  \node[accent, scale=0.75] at (-5.55,0.5) {$q\vec E$};
  \draw[gray!40!black, dashed, ->, line width=0.8pt] (-4.1,-1.15) -- (-4.1,1.15);
  \node[gray!40!black, scale=0.75, align=center] at (-3.35,0) {push\\[-2pt]$d$};
  % --- right panel: -q ---
  \draw[cyan!65!black, ->, line width=0.8pt] (2.6,1.9) -- (2.6,-1.9);
  \draw[cyan!65!black, ->, line width=0.8pt] (6.6,1.9) -- (6.6,-1.9);
  \filldraw[blue!55!black] (4.6,-1.15) circle (0.1);
  \node[blue!55!black, scale=0.85] at (4.6,-1.55) {$-q$};
  \draw[accent, ->, line width=0.9pt] (4.95,-0.95) -- (4.95,-0.05);
  \node[accent, scale=0.75] at (5.55,-0.5) {$q\vec E$};
  \draw[gray!40!black, dashed, ->, line width=0.8pt] (4.1,1.15) -- (4.1,-1.15);
  \node[gray!40!black, scale=0.75, align=center] at (3.35,0) {push\\[-2pt]$d$};
\end{tikzpicture}
\end{center}

\begin{itemize}[itemsep=9pt,parsep=4pt,topsep=4pt]
\item Force on the $+q$: $q\vec E$ points \ANS{downward, the same direction
as $\vec E$} --- that's its \emph{natural} (free) direction of motion.
\item Force on the $-q$: $q\vec E$ points \ANS{upward, opposite $\vec E$} ---
the sign flip turns the natural direction around.
\end{itemize}

Push each charge a distance $d$ \emph{against} its own natural direction ---
the same idea as raising the ball.

\[
\Delta U_+ = \ANS{qEd} \qquad\qquad \Delta U_- = \ANS{|q|Ed}
\]

\begin{itemize}[itemsep=9pt,parsep=4pt,topsep=4pt]
\item Both are \ANS{positive} --- pushing \emph{against} a charge's natural
direction always stores positive PE, regardless of the charge's sign.
\item Let either charge go instead, and it moves along its \emph{natural}
direction: PE \ANS{decreases} either way --- the field does the work.
\end{itemize}

\checkpoint{A negative charge is released in this downward field. Which way
does it move, and what happens to its PE?}

{\small\color{accent} It moves \emph{upward} --- opposite $\vec E$, since the
force on a negative charge is $q\vec E$ with $q<0$. Its PE decreases, same
rule as always: released charges move toward lower PE.}

%============================================================
\nsec{Gravitational PE, the full (non-uniform) form}

$U=mgh$ only works because $g$ barely changes over any height you can
reach near Earth's surface. Zoom out far enough to leave that ``uniform
field'' approximation, and the real, general form (valid at any separation
$r$ between two masses) is:

\[
U(r) = \ANS{-\dfrac{GMm}{r}}
\]

\begin{center}
\begin{tikzpicture}[x=1cm,y=1cm]
  \draw[gray!70, ->] (0,-6.3) -- (0,0.6) node[above, scale=0.8] {$U$};
  \draw[gray!70, ->] (-0.3,0) -- (7.6,0) node[right, scale=0.8] {$r$};
  \draw[gray!50, dashed] (0,0) -- (7.3,0);
  \draw[accent, line width=1.1pt, domain=1:7.3, samples=120, smooth]
    plot (\x, {-5.6/\x});
  \draw[gray!50] (1,0.08) -- (1,-0.08);
  \node[gray, scale=0.75] at (1,-0.4) {$R$};
  \node[gray, scale=0.7] at (-0.32,0.28) {$0$};
\end{tikzpicture}
\end{center}

\begin{itemize}[itemsep=9pt,parsep=4pt,topsep=4pt]
\item As $r\to\infty$, $U\to$ \ANS{$0$} --- same rule as always: zero PE at
infinite separation.
\item Every value of $U(r)$ here is \ANS{negative}. But moving to a
\emph{larger} $r$ (further from Earth, ``higher'') always makes $U$
\ANS{less negative --- i.e.\ larger, higher on the graph}. Don't let the
minus sign fool you: higher is \emph{always} more PE, on this curve or any
other.
\end{itemize}

\checkpoint{True or false: since $U(r)$ is always negative here, moving
\emph{further} from Earth makes $U$ smaller.}

{\small\color{accent} False. Further away makes $U$ \emph{larger} (less
negative, closer to zero) --- the graph is rising the whole way, it just
never crosses above 0. ``Negative'' and ``smaller'' are not the same thing.}

%============================================================
\nsec{Electric PE of two point charges (from the class prep)}

Pulling a pair of point charges apart, or pushing them together, takes work
--- work that gets stored as potential energy, exactly like the point-mass
form above.

\[
U(q_1,q_2) = \ANS{\dfrac{1}{4\pi\varepsilon_0}\dfrac{q_1 q_2}{r_{12}}}
\]

Sketch $U$ vs.\ $r_{12}$ below, for an opposite-sign pair and a like-sign
pair, on the same axes. (Same shape you just drew for gravity --- plus its
mirror image.)

\begin{center}
\begin{tikzpicture}[x=1cm,y=1cm]
  \draw[gray!70, ->] (0,-3.6) -- (0,3.6) node[above, scale=0.8] {$U$};
  \draw[gray!70, ->] (-0.3,0) -- (8.6,0) node[right, scale=0.8] {$r_{12}$};
  \draw[gray!50, dashed] (0.55,0) -- (7.3,0);
  \draw[red!65!black, line width=1.1pt, domain=0.55:7.3, samples=120, smooth]
    plot (\x, {1.9/\x}) node[right, scale=0.72, yshift=4pt] {like};
  \draw[blue!55!black, line width=1.1pt, domain=0.55:7.3, samples=120, smooth]
    plot (\x, {-1.9/\x}) node[right, scale=0.72, yshift=-4pt] {opposite};
\end{tikzpicture}
\end{center}

\begin{itemize}[itemsep=9pt,parsep=4pt,topsep=4pt]
\item As $r_{12}\to\infty$, $U\to$ \ANS{$0$}.
\item For three or more charges, sum over every \ANS{pair} (not every charge).
\item Opposite charges ($+$ and $-$): $U$ is \ANS{negative} at any finite
separation, and $U\to0$ only as $r\to\infty$. So $U=0$ is actually the
\ANS{largest} value here --- pulling the charges closer always finds a
\emph{lower} (more negative) PE.
\item Like charges (same sign): $U$ is \ANS{positive} at any finite separation,
and again $U\to0$ as $r\to\infty$. Here $U=0$ \emph{is} the \ANS{smallest}
value --- you can't push like charges to a lower PE than infinitely far
apart.
\end{itemize}

\[
U_{\text{total}} = \ANS{\dfrac{1}{4\pi\varepsilon_0}\!\!\sum_{\text{pairs }(i,j)}\!\dfrac{q_i q_j}{r_{ij}}}
\]

\checkpoint{Two $+q$ charges are brought from far apart to a separation $r$.
Is $U$ positive or negative? What about a $+q$ and a $-q$?}

{\small\color{accent} $+q$ and $+q$: positive (repulsive, energy stored pushing
them together). $+q$ and $-q$: negative (attractive) --- just like the ball
in the hole.}

%============================================================
\nsec{The dipole in a field}

A dipole is a $+q$ and a $-q$ separated by $\vec d$ (dipole moment
$\vec p = q\vec d$), sitting in a uniform field $\vec E$, tilted at angle
$\theta$ between $\vec p$ and $\vec E$. Each charge sits at a different
place along $\vec E$ --- use \S2's pushing idea on each charge, then add.

{\small\color{accent} Split $\vec d$ into a component along $\vec E$
($d\cos\theta$) and a component across it. Only the along-$\vec E$ part
matters (\S2): the $+q$ sits $\tfrac12 d\cos\theta$ further along $\vec E$
than center, the $-q$ sits $\tfrac12 d\cos\theta$ behind it. Adding the two
charges' PE (each $\propto -qEd$, from \S2, since moving \emph{with} $\vec
E$ lowers PE) gives $U=-qEd\cos\theta = -\vec p\cdot\vec E$.}

\[
U(\theta) = \ANS{-pE\cos\theta}
\]

\begin{itemize}[itemsep=9pt,parsep=4pt,topsep=4pt]
\item At $\theta=90^\circ$ (perpendicular to the field): $U=$ \ANS{$0$}.
\item At $\theta=0^\circ$ (aligned with the field): $U=$ \ANS{$-pE$} ---
which is \ANS{lower} than the $\theta=90^\circ$ value!
\item At $\theta=180^\circ$ (anti-aligned): $U=$ \ANS{$+pE$}.
\end{itemize}

\checkpoint{Which orientation is the most stable (lowest PE)? Is that the
orientation where $U=0$? What does this have in common with the ball in
the hole?}

{\small\color{accent} $\theta=0^\circ$ (aligned) is most stable, at
$U=-pE$ --- \emph{not} the $U=0$ orientation. Just like the ball in the
hole: $U=0$ is just wherever we started measuring from, not necessarily
the lowest energy the system can reach.}

%============================================================
\nsec{Sketch it: $U$ vs.\ $\theta$}

\begin{center}
\begin{tikzpicture}[x=1cm,y=1cm]
  \draw[gray!70, ->] (0,-2.1) -- (0,2.3) node[above, scale=0.8] {$U$};
  \draw[gray!70, ->] (0,0) -- (9.2,0) node[right, scale=0.8] {$\theta$};
  \foreach \x/\lab in {0/{$0^\circ$}, 2.2/{$90^\circ$}, 4.4/{$180^\circ$}, 6.6/{$270^\circ$}, 8.8/{$360^\circ$}}{
    \draw[gray!50] (\x,0.08) -- (\x,-0.08);
    \node[gray, scale=0.75] at (\x,-0.4) {\lab};
  }
  \draw[accent, line width=1pt, domain=0:8.8, samples=100]
    plot (\x, {-1.9*cos(\x*360/8.8)});
\end{tikzpicture}
\end{center}

{\small\color{blankink} Check your curves, and the point-charge and
uniform-field results above, against the \emph{Assembling Charges},
\emph{PE vs.\ Separation}, and \emph{Electric Dipole in a Field} sims ---
open their \textbf{Info} panels.}

\end{document}
